\chapter{Deploy e Produção}
\label{cap:deploy}

\section{Objetivo}

O professor exige preparação para ambiente de produção, mesmo com operação temporária para apresentação \cite{twelve_factor,todo.md}. Este capítulo descreve arquitetura de deploy, variáveis de ambiente e procedimentos de build.

\section{Topologia de produção proposta}

\begin{figure}[htbp]
    \centering
    \fbox{\parbox{0.92\textwidth}{\centering
    \vspace{0.4cm}
    \textbf{[Figura 4 -- Topologia de deploy]}\\[0.3cm]
    Dispositivo mobile $\rightarrow$ HTTPS $\rightarrow$ Railway (Spring Boot + PostgreSQL)\\
    Expo EAS Build $\rightarrow$ APK/AAB instalável\\
    Opcional: Vercel para landing ou painel web auxiliar\\
    \vspace{0.4cm}
    }}
    \caption{Infraestrutura de produção do Cardly}
    \label{fig:deploy-topologia}
\end{figure}

Componentes \cite{railway_docs,vercel_docs,aws_free_tier,expo_eas_build}:
\begin{itemize}
    \item \textbf{Backend}: container Java executando JAR Spring Boot.
    \item \textbf{Banco}: PostgreSQL gerenciado (Railway plugin ou AWS RDS free tier).
    \item \textbf{Mobile}: binário Android via EAS Build; Expo Go apenas em desenvolvimento.
    \item \textbf{Front web opcional}: Vercel para documentação ou página de download.
\end{itemize}

\section{Configuração por ambiente}

Princípio twelve-factor: configuração via ambiente, não via código \cite{twelve_factor}.

\subsection{Backend (Railway)}

Variáveis obrigatórias:

\begin{table}[htbp]
    \caption{Variáveis de ambiente do backend em produção}
    \label{tab:env-backend}
    \centering
    \small
    \begin{tabular}{ll}
        \toprule
        \textbf{Variável} & \textbf{Descrição} \\
        \midrule
        SPRING\_DATASOURCE\_URL & JDBC PostgreSQL \\
        SPRING\_DATASOURCE\_USERNAME & Usuário BD \\
        SPRING\_DATASOURCE\_PASSWORD & Senha BD \\
        CARDLY\_JWT\_SECRET & Segredo HS256 (256+ bits) \\
        CARDLY\_JWT\_EXPIRATION\_SECONDS & TTL token (ex.: 3600) \\
        CARDLY\_GOOGLE\_CLIENT\_ID & Client ID OAuth Google \\
        SERVER\_PORT & Porta exposta (Railway injeta) \\
        \bottomrule
    \end{tabular}
\end{table}

Flyway executa migrations no startup \cite{flyway_docs}. Health check em \texttt{/actuator/health} (se habilitado).

\subsection{Frontend mobile (Expo)}

\begin{table}[htbp]
    \caption{Variáveis públicas do app Expo}
    \label{tab:env-frontend}
    \centering
    \small
    \begin{tabular}{ll}
        \toprule
        \textbf{Variável} & \textbf{Descrição} \\
        \midrule
        EXPO\_PUBLIC\_API\_URL & URL base HTTPS da API Railway \\
        EXPO\_PUBLIC\_GOOGLE\_CLIENT\_ID & Client ID para login Google \\
        \bottomrule
    \end{tabular}
\end{table}

Rebuild obrigatório após alterar variáveis \texttt{EXPO\_PUBLIC\_*} \cite{expo_docs}.

\section{Pipeline de deploy do backend}

Procedimento resumido \cite{railway_docs}:
\begin{enumerate}
    \item Conectar repositório GitHub \texttt{cardly-backend} ao Railway.
    \item Adicionar serviço PostgreSQL e vincular \texttt{DATABASE\_URL}.
    \item Configurar variáveis da Tabela~\ref{tab:env-backend}.
    \item Deploy automático a cada push na branch principal.
    \item Validar logs: Flyway migrations OK, Tomcat started.
\end{enumerate}

Alternativa AWS Free Tier \cite{aws_free_tier}: EC2 t2.micro + RDS PostgreSQL, com maior overhead operacional.

\section{Build mobile para demonstração}

\begin{enumerate}
    \item Instalar EAS CLI: \texttt{npm install -g eas-cli}.
    \item Configurar \texttt{eas.json} com profile \texttt{preview} (APK).
    \item Executar \texttt{eas build -p android --profile preview} \cite{expo_eas_build}.
    \item Distribuir APK via link Expo ou USB na apresentação.
\end{enumerate}

Para iOS, necessário Apple Developer (opcional no escopo Android-first da equipe).

\section{CORS e rede}

Backend deve permitir origem do app. Em APIs mobile puras, CORS afeta menos chamadas nativas, mas permanece relevante para testes web e ferramentas \cite{helmet_cors_spring}. Recomenda-se lista explícita de origens, não wildcard em produção.

\section{Monitoramento e logs}

Railway exibe logs stdout/stderr. Métricas básicas: latência, taxa de 5xx, uso de memória JVM. Para demo acadêmica, monitoramento avançado (Prometheus/Grafana) é opcional.

\section{Plano de rollback}

\begin{itemize}
    \item Railway permite redeploy de release anterior.
    \item Migrations Flyway devem ser backward-compatible ou acompanhadas de scripts de reversão.
    \item APK anterior mantido localmente como fallback na apresentação.
\end{itemize}

\section{Checklist pré-apresentação}

\begin{enumerate}
    \item API pública respondendo \texttt{GET /api/me} com token válido.
    \item Banco populado com dados de demonstração (usuário admin + decks exemplo).
    \item APK compilado com \texttt{EXPO\_PUBLIC\_API\_URL} de produção.
    \item Segredo JWT de produção diferente do desenvolvimento.
    \item Google OAuth configurado com SHA-1 do keystore de release.
    \item Teste completo T01--T09 (Capítulo~\ref{cap:testes}) em rede móvel (4G/Wi-Fi).
\end{enumerate}

\section{Descomissionamento}

Conforme escopo acadêmico, após apresentação a equipe pode:
\begin{itemize}
    \item Desativar serviço Railway e banco.
    \item Revogar client IDs OAuth de teste.
    \item Manter repositórios públicos como portfólio.
\end{itemize}

\section{Production readiness -- síntese}

\begin{table}[htbp]
    \caption{Matriz de prontidão para produção}
    \label{tab:production-readiness}
    \centering
    \footnotesize
    \begin{tabular}{p{3.5cm} p{2cm} p{7cm}}
        \toprule
        \textbf{Aspecto} & \textbf{Status} & \textbf{Observação} \\
        \midrule
        Config externalizada & OK & Env vars Railway/Expo \\
        Migrations versionadas & OK & Flyway \\
        Autenticação segura & Parcial & Falta refresh token \\
        HTTPS & OK & Railway default \\
        Build reprodutível & OK & Maven + EAS \\
        Observabilidade & Básico & Logs Railway \\
        Rate limiting & Pendente & Aceitável para demo \\
        Backup BD & Manual & Snapshot Railway pré-demo \\
        \bottomrule
    \end{tabular}
\end{table}

Referências: The Twelve-Factor App \cite{twelve_factor}; documentação Railway e Expo \cite{railway_docs,expo_eas_build}.
